\documentclass[11pt]{article}
\usepackage[margin=1in]{geometry}
\usepackage{amsmath, amssymb, amsthm}

\title{Problem 466: Distinct Terms in a Multiplication Table}
\author{}
\date{}

\newtheorem{theorem}{Theorem}
\newtheorem{lemma}[theorem]{Lemma}

\begin{document}
\maketitle

\section{Problem Statement}

Let $P(m, n)$ be the number of distinct terms in an $m \times n$ multiplication table. Given: $P(64, 64) = 1263$, $P(12, 345) = 1998$, $P(32, 10^{15}) = 13826382602124302$. Find $P(64, 10^{16})$.

\section{Mathematical Foundation}

\begin{theorem}[Set-theoretic formulation]
\label{thm:set}
$P(m, n) = |\{v \in \mathbb{Z}^+ : \exists\, d \mid v,\; d \le m,\; v/d \le n\}|$.
\end{theorem}

\begin{proof}
Entry $v = i \cdot j$ appears iff $i \le m$ and $j = v/i \le n$, so $v$ appears iff it has a divisor in $[1, m]$ with complementary quotient in $[1, n]$.
\end{proof}

\begin{theorem}[Row-by-row decomposition]
\label{thm:row}
$P(m, n) = \sum_{i=1}^{m} |S_i \setminus \bigcup_{k<i} S_k|$ where $S_i = \{ij : 1 \le j \le n\}$.
\end{theorem}

\begin{proof}
Standard decomposition of a union into disjoint incremental contributions.
\end{proof}

\begin{lemma}[New element characterization]
\label{lem:new}
A value $v = ij$ in row~$i$ is new iff $i$ is the smallest divisor of~$v$ in $[1, m]$ with $v/i \le n$.
\end{lemma}

\begin{proof}
If $v$ appeared in row $k < i$, then $k \mid v$ and $v/k \le n$. Conversely, if no such~$k$ exists, $v$ is new.
\end{proof}

\begin{theorem}[M\"obius inclusion-exclusion]
\label{thm:mobius}
For each row~$i$, the count of new elements is computed by M\"obius inversion over the divisor lattice of numbers up to~$m$:
\[
\mathrm{new}_i = \sum_{T} \mu(T) \cdot \left\lfloor \frac{n}{\mathrm{lcm}(i, T)/i} \right\rfloor
\]
where the sum ranges over subsets~$T$ of ``blocking'' rows $k < i$ that share multiples with row~$i$, and $\mu(T)$ is the M\"obius coefficient for the inclusion-exclusion.
\end{theorem}

\begin{proof}
For each $j \in [1, n]$, the value $ij$ is old iff some $k < i$ divides $ij$ with $ij/k \le n$. Inclusion-exclusion over such~$k$ gives the stated formula. Each term reduces to a floor-division sum.
\end{proof}

\begin{lemma}[Hyperbola method]
\label{lem:hyp}
Sums of the form $\sum_{j=1}^{N} \lfloor N/j \rfloor$ are computable in $O(\sqrt{N})$ time.
\end{lemma}

\begin{proof}
Split at $j = \lfloor \sqrt{N} \rfloor$: compute directly for $j \le \sqrt{N}$, and group by quotient value for $j > \sqrt{N}$. Both ranges contribute $O(\sqrt{N})$ terms.
\end{proof}

\section{Algorithm}

\begin{enumerate}
    \item For each row $i = 1, \ldots, m$:
    \begin{enumerate}
        \item Identify blocking rows $k < i$ sharing common multiples with row~$i$.
        \item Apply M\"obius inclusion-exclusion to count new elements.
        \item Use the hyperbola method for efficient floor-sum evaluation.
    \end{enumerate}
    \item Sum all $\mathrm{new}_i$ to get $P(m, n)$.
\end{enumerate}

\section{Complexity Analysis}

\begin{itemize}
    \item \textbf{Time:} $O(m \cdot \sqrt{n} \cdot 2^{\omega(m)})$ where $\omega(m)$ counts distinct prime factors of numbers up to~$m$. For $m = 64$, $n = 10^{16}$: feasible.
    \item \textbf{Space:} $O(m)$.
\end{itemize}

\section{Answer}

\[
\boxed{258381958195474745}
\]

\end{document}
